\documentclass[11pt]{article}

\usepackage[margin=1in]{geometry}
\usepackage{amsmath,amssymb,bm}
\usepackage{booktabs}
\usepackage{hyperref}

\newcommand{\Rnav}{\mathcal{R}_{\mathrm{nav}}}
\newcommand{\sigmar}{\sigma_r}
\newcommand{\sigmav}{\sigma_v}
\newcommand{\GammaGeom}{\Gamma}
\newcommand{\xstate}{\bm{x}}
\newcommand{\zhat}{\hat{\bm{z}}}
\newcommand{\Phimat}{\bm{\Phi}}
\newcommand{\Qmat}{\bm{Q}}
\newcommand{\Rmat}{\bm{R}}
\newcommand{\Pmat}{\bm{P}}
\newcommand{\Kmat}{\bm{K}}
\newcommand{\Hmat}{\bm{H}}
\newcommand{\Imat}{\bm{I}}


\title{Detailed Explanation of the Navigation Risk Equation}
\author{Fabio Facin}
\date{\today}

\begin{document}
\maketitle

\begin{abstract}
This note explains each part of the proposed navigation risk score for
autonomous cislunar spacecraft navigation. The goal is twofold: first, to state
the equation in a technically precise way for aerospace and estimation
discussion; second, to give a plain-language explanation that can be understood
without advanced mathematics. The equation is not a new law of physics. It is
an engineering risk metric that tells an Extended Kalman Filter when the
spacecraft's navigation estimate should be treated with more caution.
\end{abstract}

\section{The Full Equation}

The proposed navigation risk score is
\begin{equation}
\boxed{
\Rnav(t)
=
w_r \frac{\sigmar(t)}{r_0}
+
w_v \frac{\sigmav(t)}{v_0}
+
w_g \frac{1}{\GammaGeom(t)+\epsilon}
+
w_m M(t)
+
w_l L(t)
}.
\end{equation}
It combines five kinds of warning signals:
\begin{enumerate}
  \item position uncertainty,
  \item velocity uncertainty,
  \item weak measurement geometry,
  \item missed measurements,
  \item poor lighting or visibility.
\end{enumerate}
The output $\Rnav(t)$ is dimensionless. Larger values mean that the navigation
solution is less trustworthy.

\section{Where the Equation Comes From}

The Extended Kalman Filter estimates a spacecraft state vector,
\begin{equation}
  \xstate = [x,\ y,\ \dot{x},\ \dot{y}]^T,
\end{equation}
and also maintains an uncertainty matrix,
\begin{equation}
  \Pmat =
  \begin{bmatrix}
    \Pmat_{rr} & \Pmat_{rv}\\
    \Pmat_{vr} & \Pmat_{vv}
  \end{bmatrix}.
\end{equation}
The upper-left block $\Pmat_{rr}$ describes position uncertainty, and the
lower-right block $\Pmat_{vv}$ describes velocity uncertainty. The risk
equation starts from the idea that the filter should react not only to the
current state estimate, but also to the conditions under which that estimate
was produced. A navigation estimate made with small covariance, good viewing
geometry, and consistent measurements should be treated differently from an
estimate made during sensor dropouts and poor visibility.

\section{Position-Uncertainty Term}

The first term is
\begin{equation}
  w_r \frac{\sigmar(t)}{r_0},
  \qquad
  \sigmar(t)=\sqrt{\mathrm{tr}(\Pmat_{rr}(t))}.
\end{equation}
The trace of $\Pmat_{rr}$ adds the variances in the two position directions.
Taking the square root converts variance-like units into position-like units.
Thus, $\sigmar$ behaves like a one-number summary of position uncertainty.

The scale $r_0$ makes the term dimensionless. If $r_0$ is chosen as the
position-error scale that is considered operationally meaningful, then
$\sigmar/r_0$ says how large the filter's position uncertainty is compared with
that reference level.

\subsection*{Plain-language version}

This part asks: ``How unsure are we about where the spacecraft is?'' If the
spacecraft could be in a wide region of space, this term gets bigger. If the
spacecraft's position is tightly known, this term stays small.

\section{Velocity-Uncertainty Term}

The second term is
\begin{equation}
  w_v \frac{\sigmav(t)}{v_0},
  \qquad
  \sigmav(t)=\sqrt{\mathrm{tr}(\Pmat_{vv}(t))}.
\end{equation}
Velocity matters because a small velocity error can become a large future
position error. In cislunar space, where trajectories can be dynamically
sensitive, knowing where the spacecraft is now is not enough; the estimator
must also know how confidently it understands the spacecraft's motion.

\subsection*{Plain-language version}

This part asks: ``How unsure are we about how fast and in what direction the
spacecraft is moving?'' If velocity is uncertain, the future path becomes
uncertain too.

\section{Measurement-Geometry Term}

The third term is
\begin{equation}
  w_g \frac{1}{\GammaGeom(t)+\epsilon}.
\end{equation}
The quantity $\GammaGeom(t)$ is a geometry-strength score. In estimation,
measurements are more useful when they constrain the unknown state from
informative directions. Measurements are less useful when the viewing geometry
is weak, nearly degenerate, or poorly conditioned.

The reciprocal form is intentional. When geometry is strong, $\GammaGeom$ is
larger and the risk contribution is smaller. When geometry is weak,
$\GammaGeom$ is smaller and the risk contribution grows. The small constant
$\epsilon$ prevents division by zero and sets a maximum sensitivity scale for
extremely poor geometry.

\subsection*{Plain-language version}

This part asks: ``Are the sensors looking from a useful angle?'' Imagine trying
to locate an object using only a bad viewing angle. The measurement may exist,
but it does not help as much. Bad geometry makes navigation riskier.

\section{Missed-Measurement Term}

The fourth term is
\begin{equation}
  w_m M(t).
\end{equation}
Here $M(t)$ is a missed-measurement indicator. In the simplest version,
\begin{equation}
M(t)=
\begin{cases}
1, & \text{if a scheduled measurement is unavailable},\\
0, & \text{otherwise}.
\end{cases}
\end{equation}
This term represents sensor outages, tracking interruptions, rejected
measurements, communication gaps, or any time step where the filter cannot use
new observation data.

\subsection*{Plain-language version}

This part asks: ``Did the spacecraft miss a sensor update?'' If the spacecraft
does not get new information, its estimate can drift. Missing data increases
risk.

\section{Lighting and Visibility Term}

The fifth term is
\begin{equation}
  w_l L(t).
\end{equation}
The penalty $L(t)$ describes visibility conditions. In cislunar navigation,
optical measurements can be affected by lunar lighting, eclipse, glare, target
contrast, and unfavorable Sun-target-spacecraft geometry. This term does not
try to model a full camera or star tracker. It is a controlled penalty that
lets the simulation test whether a filter reacts sensibly when optical
navigation conditions degrade.

\subsection*{Plain-language version}

This part asks: ``Can the spacecraft actually see what it needs to see?''
Bright glare, shadow, or poor contrast can make optical navigation less
reliable.

\section{The Weights}

The constants
\begin{equation}
  w_r,\quad w_v,\quad w_g,\quad w_m,\quad w_l
\end{equation}
control the importance of each warning signal. They are not physical constants.
They are tuning parameters. A serious research version of this project must
not choose them arbitrarily and then claim success. Instead, it should scan
the weights, identify stable regions, reject unstable regions, and report
failure cases.

\subsection*{Plain-language version}

The weights decide how much each warning light matters. For example, if missed
measurements are the most dangerous problem in a scenario, $w_m$ should matter
more. If bad viewing angle is the main issue, $w_g$ should matter more.

\section{How the Risk Score Changes the EKF}

The score can adapt measurement noise:
\begin{equation}
  \Rmat(t)
  =
  \Rmat_0
  \left[
  1+\beta\,\mathrm{clip}(\Rnav(t),0,R_{\max})
  \right],
\end{equation}
or process noise:
\begin{equation}
  \Qmat(t)
  =
  \Qmat_0
  \left[
  1+\alpha\,\mathrm{clip}(\Rnav(t),0,R_{\max})
  \right].
\end{equation}
Increasing $\Rmat$ tells the filter to trust measurements less. Increasing
$\Qmat$ tells the filter to admit that its dynamics prediction may be less
reliable. The clipping operation prevents the risk score from inflating the
noise matrices without limit.

\subsection*{Plain-language version}

The risk score tells the navigation computer how cautious to be. If the risk
is low, the filter behaves normally. If the risk is high, the filter becomes
more conservative about what it trusts.

\section{What Would Make the Formula Successful?}

The formula is useful only if simulations show that it improves estimator
behavior under stress. Good evidence would include:
\begin{itemize}
  \item lower RMS position error during missed measurements,
  \item reduced divergence during bad geometry,
  \item reasonable covariance growth rather than overconfidence,
  \item no major damage in clean measurement cases,
  \item clear failure boundaries when weights are too aggressive.
\end{itemize}
This is why the project uses parameter scans and controlled stress scenarios.
The goal is not to prove that the equation is always right. The goal is to
find where it helps, where it fails, and whether the idea is worth improving.

\section{One-Minute Explanation for Anyone}

A spacecraft does not truly know exactly where it is. It has a best estimate
and a measure of how uncertain that estimate is. The proposed formula is like a
navigation confidence meter. It looks at five things: how uncertain the
spacecraft is about its position, how uncertain it is about its velocity, how
good the sensor viewing angle is, whether measurements were missed, and
whether visibility is poor. If those warning signs get worse, the risk score
goes up. Then the Kalman filter can become more cautious instead of blindly
trusting bad information.

\section{Short Professor-Facing Summary}

This project proposes an interpretable risk metric for adaptive EKF tuning in
a planar CR3BP cislunar navigation simulation. The contribution is not a new
dynamics model and not a new physical law. The contribution is a testable
estimator architecture: use covariance, geometry, dropout, and visibility
signals to adapt process or measurement noise, then compare against a fixed
EKF under controlled stress cases.

\end{document}
